% Ergebnisse: 1-1 1/2 Seiten (Folie 16). Nur Werte aus results/ms1/REPORT.md, Mittelwert +- Std über 3 Seeds.
\section{Ergebnisse}
\label{sec:results}

\paragraph{Vergleich aller Methoden.}
\Cref{tab:methods} und \cref{fig:efficiency} zeigen alle Methoden. LoRA mit $r=8$ erreicht die
höchste Test-Accuracy ($96{,}21\pm0{,}24$\,\%), gefolgt von LoRA $r=4$ ($96{,}11\pm0{,}21$\,\%) und
Linear Probing ($96{,}03\pm0{,}12$\,\%), das nur 0,02\,\% der Parameter trainiert. Die Unterschiede
zwischen diesen drei liegen innerhalb der Seed-Streuung. Volles Fine-Tuning ($95{,}39\pm0{,}39$\,\%)
liegt 0,82 Prozentpunkte (pp) unter der besten LoRA-Variante, mehr als unsere Schwelle. Die CNNs
erreichen $94{,}78\pm0{,}46$\,\% (ConvNeXt-T) und $94{,}22\pm0{,}85$\,\% (ResNet-50); die beste
ViT-Methode liegt 1,42\,pp über dem besten CNN. Bei der Inferenz braucht der ViT-B/16 jedoch 35,1\,GFLOPs
pro Bild, ResNet-50 8,2 und ConvNeXt-T 8,9.

\begin{table*}[t]
  \centering
  \caption{Testergebnisse (Mittelwert $\pm$ Standardabweichung über 3 Seeds; fett: bester Wert).
  Trainierbar: im Training aktualisierte Parameter. Zeit: Training eines Laufs auf einer RTX~5090. GFLOPs: Rechenaufwand
  pro Bild bei der Inferenz. LoRA-Änderungen werden nach dem Training eingerechnet (keine zusätzlichen
  Inferenzkosten).}
  \label{tab:methods}
  \small
  \begin{tabular}{lccccc}
    \toprule
    Methode & Acc.\ (\%) & Macro-F1 (\%) & Trainierbar & Zeit (min) & GFLOPs \\
    \midrule
    \methodrows
    \bottomrule
  \end{tabular}
\end{table*}

\begin{figure}[t]
  \centering
  \resultfigure{paper/effizienz.png}{\linewidth}{Accuracy über trainierbaren Parametern}
  \caption{Test-Accuracy über der Zahl trainierbarer Parameter (logarithmische Skala). Punkte: Mittelwert über
  3 Seeds, Balken: $\pm$ eine Standardabweichung. Mit weniger trainierbaren Parametern bleibt die Accuracy
  gleich oder höher; volles Fine-Tuning und die CNNs liegen darunter.}
  \label{fig:efficiency}
\end{figure}

\paragraph{LoRA-Parameter.}
Mit Anpassung von Query und Value beträgt die Accuracy 96,11, 96,21 und 95,96\,\% für $r=4$, 8 und 16
(0,17, 0,31 und 0,61\,M trainierbare Parameter). Die Anpassung aller linearen Layer mit $r=8$
(1,35\,M Parameter) ergibt $95{,}82\pm0{,}18$\,\%, 0,39\,pp unter Query/Value allein und damit innerhalb
der Seed-Streuung.

\paragraph{Trainingsverlauf.}
Die beste Validation-Accuracy erreicht Linear Probing nach $3{,}7\pm1{,}2$ Epochen, die übrigen Methoden
nach 5--7 (\cref{fig:curves}).

\begin{figure}[h]
  \centering
  \resultfigure{paper/lernkurven.png}{\linewidth}{Lernkurven}
  \caption{Validation-Accuracy je Epoche (Mittelwert über 3 Seeds; die y-Achse beginnt bei 68\,\%). Linear
  Probing liegt schon nach einer Epoche bei 94\,\%. Die CNNs starten niedriger (ResNet-50: 70\,\%) und
  bleiben nach dem Aufholen unter den LoRA-Varianten.}
  \label{fig:curves}
\end{figure}

\paragraph{Augmentierung.}
Für LoRA $r=8$ ergibt \emph{light} $96{,}21\pm0{,}24$\,\%, \emph{none} $95{,}99\pm0{,}16$\,\% und
\emph{strong} $95{,}91\pm0{,}36$\,\%; die Unterschiede liegen innerhalb der Seed-Streuung.

\paragraph{Fehler.}
Für LoRA $r=4$ (auf Validation gewählt) sind American Pit Bull Terrier (F1 77,2\,\%), Staffordshire
Bull Terrier (78,2\,\%) und American Bulldog (88,7\,\%) am schwierigsten (Anhang~\ref{sec:perbreed}).
Häufigste Verwechslung: Pit Bull $\rightarrow$ Staffordshire (62 Bilder).
